\chapter{Literature Review}
\label{chap:litrev}

Reading the literature on retrieval-augmented generation, autonomous agents,
and multimodal learning is, at this point, a somewhat daunting exercise---the
field has moved fast enough that survey papers published in mid-2023 describe a
landscape that already looks somewhat dated compared to what is in production
today. What follows is not an attempt to be exhaustive; it is an effort to trace
the specific ideas that directly shaped the design decisions in this project.

\section{The RAG Lineage}
\label{sec:rag-lineage}

The formal starting point is \citet{lewis2020rag}, who showed that pairing a
parametric sequence-to-sequence model with a non-parametric dense retriever over
Wikipedia documents substantially improved performance on open-domain question
answering without additional fine-tuning. The key insight---that a model can look
things up rather than memorise them---sounds obvious in retrospect, but the
clean experimental demonstration of it was genuinely useful for the field.

What that paper does not address is what to do when retrieved documents are
irrelevant, contradictory, or simply unhelpful. That is the problem that
Corrective RAG (CRAG) was designed to tackle. \citet{yan2024crag} introduce an
automatic evaluator that sits between retrieval and generation and redirects the
system to a web search when the initial vector retrieval scores poorly on
relevance. We implemented a version of this in our \texttt{CRAGPipeline} class;
the evaluator runs as part of the single-hop retrieval flow and is one of the
more practically useful components we adapted from the literature.

\citet{gao2023retrieval} provide the broadest survey of the space, organising
the field into Naive RAG (retrieve, stuff into context, generate), Advanced
RAG (add query rewriting, re-ranking, and compression), and Modular RAG
(swappable components). The OmniRAG pipeline sits firmly in the Modular RAG
category: query rewriting, HyDE, hybrid search, FlashRank neural re-ranking,
CRAG, contextual compression, two-stage generation, self-critique, and quality
scoring are all independent modules wired together by a single \texttt{process\_query}
controller.

\subsection{HyDE: Searching for Answers Rather Than Questions}

One of the more clever retrieval tricks we encountered---and adopted---is
Hypothetical Document Embeddings. \citet{gao2022precise} observe that when you
embed a question for dense retrieval, you are embedding a linguistic object that
looks nothing like the documents you are trying to retrieve, which are answers
rather than questions. The fix is to ask the LLM to generate a plausible
hypothetical answer to the question, embed \textit{that}, and use the resulting
vector for retrieval. In technical domains, where questions use different
vocabulary from the explanations that answer them, this meaningfully improves
retrieval recall. The effect is particularly visible for questions that use
colloquial phrasing to ask about domain-specific concepts.

\subsection{Graph RAG: Communities as Retrieval Units}

Dense retrieval over individual text chunks has a well-known weakness: it cannot
easily answer questions that require connecting information from multiple
disparate documents. \citet{edge2024local} address this with GraphRAG, which
builds a knowledge graph over the corpus and retrieves at the level of community
summaries produced by graph partitioning (specifically, the Louvain algorithm).
For multi-hop queries---``compare how X approaches the problem in document A
versus how Y handles it in document B''---this is a considerably more effective
retrieval unit than a 512-token text chunk. Our \texttt{KnowledgeGraph}
service implements this using the \texttt{python-louvain} library and stores
community summaries in MongoDB for low-latency access during the multi-hop
retrieval path.

\section{Autonomous Research Agents}
\label{sec:agents-review}

The closest conceptual ancestor of our \texttt{DeepResearchAgent} is
ReAct~\citep{yao2022react}, which interleaves language model reasoning steps
(``thoughts'') with tool invocations (``actions'') and observations in a single
context. The loop---think, act, observe, think again---is surprisingly effective
for tasks that benefit from dynamic, evidence-driven reasoning chains.

AutoGen~\citep{wu2023autogen} extends this to multi-agent conversations where
different agents adopt distinct roles (orchestrator, coder, reviewer) and
collaborate through structured message passing. Our research agent is
single-agent but draws on the same principle: distinct numbered phases (Decompose,
Search, Evaluate, Refine, Synthesize) each have a defined input contract and
output contract, which makes the pipeline controllable and its failure modes
localised to individual phases rather than opaque across the whole workflow.

The iterative gap-analysis step---where the agent inspects which sub-questions
remain under-sourced and generates targeted follow-up queries---was our own
addition. It was motivated by observing, during early testing, that a single
research pass reliably missed nuances in broad comparative questions. The
repetition of the Search$\to$Evaluate loop until a minimum source count is
reached, or a maximum iteration budget is exhausted, produced meaningfully more
comprehensive outputs in practice.

Reflexion~\citep{shinn2023reflexion} contributes the idea of verbal self-
evaluation: rather than discarding failed actions, the model generates a natural-
language critique of what went wrong and uses that critique to guide the next
attempt. Our \texttt{SelfCritique} module applies this at the generation level:
the model scores its own answer against the retrieved documents and emits a
confidence value that propagates to the frontend as response metadata.

\section{Multimodal Perception}
\label{sec:multimodal-review}

CLIP~\citep{radford2021clip} is the backbone of our visual understanding layer.
Its contrastive pre-training on 400 million image-text pairs from the web
produces a joint embedding space where images and their textual descriptions map
to nearby points. For our purposes this means that image search and text search
can share the same FAISS vector index: a user can ask a question in text and
retrieve relevant images, or upload an image and retrieve relevant text passages,
without any structural changes to the retrieval pipeline.

We supplement CLIP with EasyOCR (for extracting printed and handwritten text
from images) and YOLOv8~\citep{jocher2023yolov8} (for detecting and labelling
objects in images). The three-component stack covers the main categories of
engineering images that students actually upload: photos of handwritten derivations
(OCR-dominant), circuit diagrams and block diagrams (OCR + YOLO), and
oscilloscope or simulator screenshots (YOLO + CLIP context).

\section{Secure Code Execution}
\label{sec:coding-review}

Online judge systems have grappled with safe code execution for decades. The
standard approach---execute submitted code inside resource-limited containers
with network isolation disabled---is well-established in competitive programming
infrastructure. Our \texttt{SecureSandbox} follows roughly the same model:
ephemeral Docker containers with CPU and memory cgroups, wall-time enforcement,
and stdin/stdout redirection through a controlled channel. What differs from a
typical OJ setup is the addition of a WebSocket terminal (via XTerm.js) that
gives the student an interactive shell rather than just batch execution output,
and a DAP-compliant debugger that supports breakpoints and variable inspection
from within the browser.

\section{Text Quality in Generated Content}
\label{sec:quality-review}

One concern that kept surfacing during early user testing was that the first-draft
outputs of the RAG pipeline, while factually grounded, were recognisably
``AI-written'' in style---heavy on filler phrases, vague hedges, and repetitive
sentence structures. \citet{guo2023llm_quality} document this phenomenon
empirically, showing that LLM-generated text is distinguishable from human-written
text along several surface-level stylistic dimensions. We designed a 10-module
Text Quality Upgrade stack (documented in \texttt{docs/architecture/}) to address
this: a \texttt{LanguageOptimizer} that detects and removes high-frequency
LLM-typical phrases, a \texttt{DensityController} that enforces an information-
to-word-count ratio floor, and an \texttt{AnswerRefiner} that iteratively
compresses wordy drafts without discarding meaning.

\section{Engineering Education Technology}
\label{sec:edu-review}

Work on technology-enhanced undergraduate education spans three decades and a
wide range of intervention types---intelligent tutoring systems, automated
formative assessment, collaborative design environments, flipped-classroom tools.
Within the AI-assisted category, the systems most relevant to this project are
the ones that blur the boundary between a study assistant and a research tool.
\citet{liu2023adaptive} review adaptive tutoring approaches that adjust difficulty
and scaffolding to individual learner models. What those systems share with
EngUnity is the idea that a single, consistent user model should persist across
multiple interactions. Where EngUnity diverges is in scope: instead of adapting
a pre-specified curriculum, it adapts the retrieval and reasoning strategy to
the query at hand, which is a considerably more open-ended problem.

\section{Positioning EngUnity}
\label{sec:positioning}

Table~\ref{tab:comparison} maps the capabilities described above against the
most commonly used tools in the target student population.

\begin{table}[ht]
    \centering
    \caption{Capability comparison across AI-augmented tools and platforms.}
    \label{tab:comparison}
    \small
    \begin{tabular}{lccccc}
        \toprule
        \textbf{Platform} & \textbf{RAG} & \textbf{Research Agent} &
        \textbf{Multimodal} & \textbf{Code Lab} & \textbf{Ed.~Focus} \\
        \midrule
        ChatGPT             & Partial & --       & \checkmark & Partial & --          \\
        Gemini Advanced     & \checkmark & --     & \checkmark & --      & --          \\
        Perplexity AI       & \checkmark & Partial & --         & --      & --          \\
        GitHub Copilot      & --      & --       & --         & \checkmark & --       \\
        Cursor              & --      & --       & --         & \checkmark & --       \\
        LeetCode            & --      & --       & --         & \checkmark & \checkmark \\
        \textbf{EngUnity}   & \textbf{\checkmark} & \textbf{\checkmark} &
        \textbf{\checkmark} & \textbf{\checkmark} & \textbf{\checkmark} \\
        \bottomrule
    \end{tabular}
\end{table}

None of the commercial tools occupies the same column profile as EngUnity.
Whether that combination is the right one for the target population is partly
an empirical question and partly a matter of what the student actually needs
in a given study session. The design deliberately makes each subsystem usable
independently so that a student who only wants the research agent, for instance,
does not need to engage with the analytics workbench or the code lab.
